\section{Detailed paper-by-paper critical appraisals}
\label{app:appraisals}
The following structured records are part of the review, not a substitute for its
thematic synthesis. Each separates a method's contribution, the inspected evidence,
critical interpretation and implications for the project. ``Selected full text''
means the identified passages were inspected; it does not mean code replication,
verification of every appendix, or equivalence between a preprint and its later
version of record. Record-limited entries remain explicit evidence limitations.

\subsection{Foundations and social requirements}
\subsubsection{From obstacle avoidance to human-aware navigation}\label{paper:survey2013}
\textbf{Publication:} \citet{survey2013}. \emph{Human-aware robot navigation: A survey}. \href{https://doi.org/10.1016/j.robot.2013.05.007}{Publication record}.

\textbf{Evidence anchor:} Selected full-text inspection. Author manuscript: introduction, human-awareness requirements, Section 3 architecture discussion. \href{https://hal.science/hal-01684295v1/document}{Inspected source copy}.

\paragraph{Method and contribution.} Kruse and colleagues organize navigation around requirements arising from people, rather than treating a person as another occupied cell. Their discussion connects comfort, socially appropriate motion and task-specific interaction to perception, prediction, local planning and deliberate behavior. The architecture distinguishes reactive responses from decisions such as following, guiding, approaching and resolving a blocked passage. This is an important conceptual foundation: social navigation is a coupled perception-and-behavior problem, not a particular local planner.

\paragraph{What the evidence establishes.} The evidence is a literature synthesis and architectural classification. The inspected discussion lists concrete behaviors such as reducing speed near a person, avoiding erratic motion and approaching appropriately for an intended interaction. It does not provide a common experimental rerun of all surveyed algorithms.

\paragraph{Critical appraisal.} The breadth is useful for defining requirements, but a taxonomy is not a validated universal utility function. Different behaviors can conflict: a direct approach may be appropriate for interaction and intrusive during independent passage. Its age also means that contemporary end-to-end humanoid results must be assessed separately. The review should not be used to characterize the capabilities of all later learning methods.

\paragraph{Implication for novel work.} Use this paper to define the task before selecting metrics. Geoff-Sim currently addresses independent navigation and experimental delivery; it should not claim guiding, negotiation or perceived comfort merely because it avoids a moving actor.

\subsubsection{A multifaceted taxonomy of socially aware navigation}\label{paper:survey2024}
\textbf{Publication:} \citet{survey2024}. \emph{A survey on socially aware robot navigation: Taxonomy and future challenges}. \href{https://doi.org/10.1177/02783649241230562}{Publication record}.

\textbf{Evidence anchor:} Selected full-text inspection. Accepted manuscript: abstract and taxonomy definitions, especially planning/decision-making, situation awareness and evaluation diagrams. \href{https://publications.aston.ac.uk/id/eprint/46055/1/Singamaneni_et_al_AAM_2311.06922_Survey_on_Socially_Aware_Robot_Navigation_taxonomy_and_future_challenges.pdf}{Inspected source copy}.

\paragraph{Method and contribution.} Singamaneni and colleagues replace a single algorithm hierarchy with several complementary taxonomies. The inspected definitions distinguish implicit and explicit communication, independent/assistive/collaborative tasks, negotiation, local motion generation and global decision-making. Situation awareness includes representing agents, predicting their evolution and interpreting environmental context. Evaluation is treated as a separate dimension rather than inferred from the planner's architecture.

\paragraph{What the evidence establishes.} The review maps a broad preceding decade of work and identifies relationships among representations, algorithms, tasks and evaluation tools. Its principal output is a structured description of the field, not a numerical leaderboard. The faceted view permits two systems using different mathematical methods to be compared on the same task or information requirement.

\paragraph{Critical appraisal.} Taxonomic coverage does not establish that any combination of components delivers reliable autonomy. Labels such as learning-based or social remain too coarse to determine deployment privilege, latency or physical contact fidelity. A framework may cover several categories while still omitting the interactions between them. This makes the taxonomy a design and reporting instrument, not a substitute for source-level comparison.

\paragraph{Implication for novel work.} Organize Geoff-Sim's contribution across task, information, embodiment and evidence dimensions. Its most promising novelty lies at an interface between those dimensions, not in introducing another instance of an already populated algorithm family.

\subsubsection{Learning-based social navigation and the role of input contracts}\label{paper:survey2023}
\textbf{Publication:} \citet{survey2023}. \emph{Social robot navigation: a review and benchmarking of learning-based methods}. \href{https://doi.org/10.3389/frobt.2025.1658643}{Publication record}.

\textbf{Evidence anchor:} Selected full-text inspection. 2025 publisher PDF: Tables 1, 2 and 4; Sections 2.4 and 2.5 on prediction, hybrid planning and safety-aware navigation. The stable citation key does not denote its publication year. \href{https://public-pages-files-2025.frontiersin.org/journals/robotics-and-ai/articles/10.3389/frobt.2025.1658643/pdf}{Inspected source copy}.

\paragraph{Method and contribution.} Alyassi and colleagues classify learning methods by what information they consume and how they make decisions. The tables separate raw-sensor end-to-end methods, human-state methods, attention-based interaction models and prediction-based planners. The discussion covers hybrid combinations in which learned subgoals or pedestrian forecasts are passed to model-predictive control, as well as policy switching and fallback mechanisms.

\paragraph{What the evidence establishes.} The inspected tables explicitly record architecture, training environment, output interface, code availability and real-world demonstrations. The prediction discussion reports that sophisticated forecasts do not uniformly dominate constant-velocity alternatives in every downstream application. Its benchmarking context is useful, but this appraisal does not extract a universal ranking or pool its outcomes with Geoff-Sim trials.

\paragraph{Critical appraisal.} The strongest lesson is that architectural labels conceal different information contracts. A detector-based method and a state-vector policy solve different parts of the pipeline. Secondary descriptions of a method are useful for discovery but do not replace reading its primary experiment. Hybrid planning, uncertainty handling and fallback switching are already established directions; combining them generically is not itself a novelty claim.

\paragraph{Implication for novel work.} Use this review to select matched baselines and identify additional primary literature before submission, particularly robust MPC and uncertainty-aware hybrid policies. Treat proposed gaps in this document as conditional on the bounded primary corpus.

\subsubsection{Social-navigation principles and evaluation design}\label{paper:guidelines}
\textbf{Publication:} \citet{guidelines}. \emph{Principles and Guidelines for Evaluating Social Robot Navigation Algorithms}. \href{https://doi.org/10.1145/3700599}{Publication record}.

\textbf{Evidence anchor:} Selected full-text inspection. HAL manuscript: principles/guidelines overview and scenario-design methodology, including scenario cards. \href{https://hal.science/hal-04164243/document}{Inspected source copy}.

\paragraph{Method and contribution.} Francis and colleagues distinguish safety, comfort, legibility, politeness, social norms/competency, understanding of agents, proactivity and contextual appropriateness. They connect these principles to experimental design, metrics, scenarios, simulators, benchmarks and datasets. Scenario cards specify research context, intended robot task, intended human behavior and success criteria, making the experimental question explicit before a metric is calculated.

\paragraph{What the evidence establishes.} This is a community methodological synthesis rather than a new robot controller. The inspected scenario-design section explains why scenarios should be representative and flexible, instead of prescribing the intended solution so narrowly that an algorithm can succeed without general competence. It also separates the quality of an evaluation instrument from the quantity of measurements collected.

\paragraph{Critical appraisal.} A principle list does not supply a single scalar social score. Some dimensions require human judgments or a context-specific instrument; others can be measured geometrically. Counting contacts does not measure politeness, and rating a short video does not establish physical reliability. The accessible manuscript predates the final bibliographic record, so version-specific wording must remain identifiable.

\paragraph{Implication for novel work.} Geoff-Sim should publish scenario cards, source contracts, failure categories and separate outcome families. This is a methodological baseline to adopt, not a new contribution to claim. Novelty must come from the scientific question or new validated mechanism tested with it.

\subsubsection{Proxemics as a contextual variable rather than a universal radius}\label{paper:proxemics2026}
\textbf{Publication:} \citet{proxemics2026}. \emph{Advancing a taxonomy of proxemics for socially aware robot navigation}. \href{https://doi.org/10.3389/frobt.2026.1800762}{Publication record}.

\textbf{Evidence anchor:} Selected full-text inspection. Publisher PDF: Section 3 methodology, Table 1 shape-model comparison and variable-classification discussion. \href{https://www.frontiersin.org/journals/robotics-and-ai/articles/10.3389/frobt.2026.1800762/pdf}{Inspected source copy}.

\paragraph{Method and contribution.} Nahum and colleagues structure proxemic influences into human, robot, environment and context factors, alongside spatial outcome variables. Their method combines a top-down taxonomy with bottom-up extraction from experimental studies. The inspected table contrasts concentric circles, asymmetric and motion-dependent shapes, and emerging three-dimensional representations. Personal space is therefore modeled as contingent on the interaction, not as a fixed property of an isolated person.

\paragraph{What the evidence establishes.} The paper describes a 2020-2025 search yielding more than 190 initial articles, 99 screened articles and 43 retained items, comprising four surveys and 39 experimental studies. These are the source paper's counts, not the screening counts of this review. Variables were manually coded, with generative-AI assistance for clustering, and underrepresented combinations were identified through frequency analysis.

\paragraph{Critical appraisal.} An infrequently studied factor is not automatically a practically important research gap. The taxonomy and corpus-selection process determine which absences become visible. Context-specific findings also resist direct numerical aggregation into one clearance threshold. A robot's body size, speed, purpose and approach direction may change judgments even when center distance is identical.

\paragraph{Implication for novel work.} Keep a fixed proxemic threshold as an explicitly declared engineering metric, then test contextual alternatives. A future humanoid-comfort contribution needs a study connecting motion, morphology and participant judgments, rather than relabeling a distance penalty as validated comfort.

\subsection{Geometric avoidance, interaction and safety}
\subsubsection{DWA and the distinction between reachable and stoppable commands}\label{paper:dwa}
\textbf{Publication:} \citet{dwa}. \emph{The dynamic window approach to collision avoidance}. \href{https://doi.org/10.1109/100.580977}{Publication record}.

\textbf{Evidence anchor:} Selected full-text inspection. Author technical manuscript: abstract, Sections 1-4, motion-model and dynamic-window explanation. \href{https://publications.ri.cmu.edu/storage/publications/pub_files/pub1/fox_dieter_1997_1/fox_dieter_1997_1.pdf}{Inspected source copy}.

\paragraph{Method and contribution.} Fox, Burgard and Thrun search directly in translational/rotational velocity space. The admissible set contains commands whose trajectories allow stopping before an obstacle, while the dynamic window restricts candidates to velocities reachable within the next control interval. The objective balances goal alignment, speed and obstacle clearance. Piecewise circular trajectories simplify collision checking for the synchro-drive model.

\paragraph{What the evidence establishes.} The manuscript describes implementation on RHINO and other mobile robots, with reported operation up to 95 cm/s in populated environments. This is historical hardware evidence under the paper's platform and sensing conditions. It is not a contemporary multi-seed humanoid crowd benchmark, and its demonstrations do not define the exposure denominator needed for a reliability estimate today.

\paragraph{Critical appraisal.} DWA directly contradicts a claim that classical avoidance ignores dynamics or braking. Its model is nevertheless a specific abstraction: planar velocity control and local trajectories do not automatically capture a learned biped policy's deadbands, asymmetry, balance transients or upper-body sweep. Local optimality and incomplete environmental information also leave global liveness problems unresolved.

\paragraph{Implication for novel work.} The relevant extension is calibrated policy-level reachability under sensing age, compared against a properly tuned admissible-velocity baseline. Adding an acceleration cap to an otherwise instantaneous planner is an engineering correction, not sufficient methodological novelty.

\subsubsection{Velocity obstacles as a relative-motion foundation}\label{paper:vo}
\textbf{Publication:} \citet{vo}. \emph{Motion Planning in Dynamic Environments Using Velocity Obstacles}. \href{https://doi.org/10.1177/027836499801700706}{Publication record}.

\textbf{Evidence anchor:} Record-limited inspection. Verified DOI metadata and recorded abstract; attempted institutional full-text locator was unavailable.

\paragraph{Method and contribution.} The velocity-obstacle formulation represents relative velocities that lead to collision with a moving obstacle and reasons about selecting motion outside that region. The retained record places dynamic-environment motion planning and velocity-space reasoning at the foundation of later reciprocal methods. The original work also discusses admissibility; the conceptual account here is deliberately separated from an uninspected reconstruction of its implementation.

\paragraph{What the evidence establishes.} The primary bibliographic identity is verified, but the full paper was not retrieved in this expansion. Consequently, no new table values, hardware protocol, runtime number or completeness claim is attributed to this publication. Detailed executable-model comparisons should use the accessible ORCA and DWA texts and return to this original article when access is obtained.

\paragraph{Critical appraisal.} The conceptual strength is interpretability: a rejection can be related to relative motion and geometry. Transfer depends on the validity of state estimates, prediction assumptions and command realizability. Those are comparison questions, not evidence that the original paper neglects every such issue. Missing access must not be converted into a negative assessment of its methods.

\paragraph{Implication for novel work.} Use VO as a foundational baseline family and novelty boundary. Resolve the original full text before making a first-of-kind claim about velocity admissibility, dynamic obstacles or uncertainty treatment.

\subsubsection{ORCA: reciprocal guarantees and their execution assumptions}\label{paper:orca}
\textbf{Publication:} \citet{orca}. \emph{Reciprocal n-Body Collision Avoidance}. \href{https://doi.org/10.1007/978-3-642-19457-3_1}{Publication record}.

\textbf{Evidence anchor:} Selected full-text inspection. Author chapter: introduction, reciprocal formulation, Section 6 experiments and Section 7 limitations. \href{https://gamma.cs.unc.edu/ORCA/publications/ORCA.pdf}{Inspected source copy}.

\paragraph{Method and contribution.} ORCA derives pairwise sufficient conditions for reciprocal collision avoidance, assigns each cooperating agent part of the required velocity correction and intersects the permitted half-planes. A low-dimensional linear program chooses a velocity close to the preferred one. The method separates local collision avoidance from the global route that provides that preference and is computationally attractive for large numbers of independently deciding agents.

\paragraph{What the evidence establishes.} The experiments include circle exchange and office evacuation with thousands of simulated agents. Reported timings distinguish scenarios and hardware: for example, 5,000 agents required 8 ms in the circle case and 15.6 ms in the office case on eight cores. These are simulator-wide algorithm timings, not observation-to-actuation latency on a walking robot.

\paragraph{Critical appraisal.} The conclusion explicitly states that its simple robot model ignores kinematic and dynamic constraints and that clamping a computed velocity afterward does not preserve strict collision avoidance. Reciprocity is equally important: nonreactive people do not implement their assigned correction. Thus performance under unilateral avoidance cannot be inferred from the reciprocal guarantee, and a baseline should disclose responsibility weights and post-processing.

\paragraph{Implication for novel work.} Compare reciprocal and nonreciprocal actor regimes separately. Adaptation to the G1 must model what commands can actually be executed, while preserving a clearly documented baseline rather than silently changing ORCA's assumptions.

\subsubsection{Sparse timed-trajectory optimization}\label{paper:teb}
\textbf{Publication:} \citet{teb}. \emph{Efficient trajectory optimization using a sparse model}. \href{https://doi.org/10.1109/ecmr.2013.6698833}{Publication record}.

\textbf{Evidence anchor:} Record-limited inspection. Verified DOI record and the authors' maintained teb\_local\_planner description; the original 2013 paper PDF was not retrieved.

\paragraph{Method and contribution.} The timed-elastic-band family jointly adjusts local trajectory geometry and timing to trade execution time, obstacle separation and kinodynamic constraints. The authors' software description identifies this optimization structure and cites the 2013 sparse-model paper. This review uses that implementation-level description as context and does not assume that every feature of the later software was already present in the cited conference paper.

\paragraph{What the evidence establishes.} The paper identity and its place in the method lineage are established. The current software repository is accessible, but source inspection and a successful software build are not experimental replication of the 2013 results. No uninspected original runtime, success rate or dynamic-obstacle guarantee is reported here.

\paragraph{Critical appraisal.} An optimized trajectory is only useful if the downstream controller follows its timing and geometry. The important comparative issues are initialization, obstacle forecasts, constraint handling and execution mismatch. These are reasons to include timed optimization as a serious baseline, not proof that a simpler response-envelope method is superior. Later homotopy and topology extensions must be cited separately if evaluated.

\paragraph{Implication for novel work.} Use an explicitly versioned TEB implementation for comparison, state its command model and shared supervisor, and obtain the original primary text before assigning paper-specific limitations or claiming a new time-aware planning formulation.

\subsubsection{Social forces and emergent pedestrian organization}\label{paper:socialforce}
\textbf{Publication:} \citet{socialforce}. \emph{Social force model for pedestrian dynamics}. \href{https://doi.org/10.1103/physreve.51.4282}{Publication record}.

\textbf{Evidence anchor:} Selected full-text inspection. Author preprint: abstract, force-model exposition and collective-behavior illustrations, including lane formation. \href{https://arxiv.org/pdf/cond-mat/9805244}{Inspected source copy}.

\paragraph{Method and contribution.} Helbing and Molnar describe pedestrian movement using an acceleration toward desired velocity, repulsive interactions with people and boundaries, and possible attractions. The forces represent motivations rather than literal mechanical contact forces. Coupled stochastic equations produce collective effects such as directional lanes, showing how locally specified interactions can generate structured crowd motion without centrally planning every pedestrian.

\paragraph{What the evidence establishes.} The inspected illustrations demonstrate emergent crowd organization and compare the modeled effects with observed pedestrian phenomena. This is evidence about a behavioral modeling mechanism, not a robot's ability to recognize social context or a universal empirical model of all people. The accessible arXiv deposit is later than the original journal publication.

\paragraph{Critical appraisal.} Parameters, boundary conditions and desired directions strongly influence the resulting crowd. Using the same social-force model both to generate training data and to evaluate a learned robot can reward adaptation to the simulator's assumptions. Apparent smoothness or lane formation does not establish realistic responses to a humanoid, vulnerable road users, distraction or cultural conventions.

\paragraph{Implication for novel work.} Treat crowd-generation policy as an experimental factor. Geoff-Sim's independent actors answer a different question from mutually adapting agents, and results should be repeated under more than one behavior model before claiming broad social competence.

\subsubsection{Cooperation and the freezing-robot problem}\label{paper:freezing}
\textbf{Publication:} \citet{freezing}. \emph{Robot navigation in dense human crowds: the case for cooperation}. \href{https://doi.org/10.1109/icra.2013.6630866}{Publication record}.

\textbf{Evidence anchor:} Record-limited inspection. Verified DOI record and Caltech institutional abstract; the landing record was metadata-only.

\paragraph{Method and contribution.} Trautman and colleagues frame dense-crowd navigation as joint behavior rather than avoidance of independently predicted obstacles. The institutional abstract describes an extension of interacting Gaussian processes to multiple goals and stochastic movement duration. Cooperation is intended to reduce the excessive conservatism that arises when the robot considers other agents' possible futures without accounting for mutual adaptation.

\paragraph{What the evidence establishes.} The accessible institutional abstract reports 488 completed runs and comparisons with teleoperation and noncooperative/reactive planners. These are abstract-reported results, not independently inspected protocols. This appraisal therefore avoids reconstructing participant allocation, unsafe-event adjudication or uncertainty intervals from that count.

\paragraph{Critical appraisal.} The central insight is relevant even when the full protocol is unavailable: treating all possible occupancy as fixed exclusion can eliminate every route, while assuming cooperation can fail when people do not yield. A model of interaction must be evaluated in conditions where its behavioral assumptions are both satisfied and violated. Access limitations prevent a detailed original-table audit here.

\paragraph{Implication for novel work.} Measure liveness explicitly: time blocked, progress after yielding and interventions, alongside contacts. A new G1 supervisor must not be presented as the first treatment of freezing or cooperative navigation; its novelty would need a specific execution- and observation-aware mechanism.

\subsubsection{Biped viability, uncertainty and relaxed passive safety}\label{paper:biped2016}
\textbf{Publication:} \citet{biped2016}. \emph{Safe navigation strategies for a biped robot walking in a crowd}. \href{https://doi.org/10.1109/humanoids.2016.7803304}{Publication record}.

\textbf{Evidence anchor:} Selected full-text inspection. HAL manuscript: Sections II-V; equations 4-12; hierarchy definitions; Tables I-V and failure discussion. \href{https://inria.hal.science/hal-01418340/document}{Inspected source copy}.

\paragraph{Method and contribution.} Bohorquez and colleagues combine balance and collision avoidance using a linearized center-of-mass/center-of-pressure model, foot-placement constraints and a terminal capture/stopping condition. Lexicographic MPC gives physical feasibility and near-term collision avoidance priority over velocity tracking. Pedestrians move at constant velocity and do not avoid the robot; position and velocity uncertainty expand their occupied region over the horizon. Alternative hierarchies defer or relax the terminal stop requirement when strict passive safety becomes counterproductive.

\paragraph{What the evidence establishes.} The numerical study uses an HRP-2-like model, 100 stored crowd realizations per condition, varying crowd size/speed, field of view, horizon and uncertainty. It explicitly excludes perception occlusions. The analysis relates observable range to closing speed and horizon, and distinguishes robot feasibility failures from collisions. Relaxing physical feasibility can reduce collisions by allowing the robot to fall, which is not an overall improvement.

\paragraph{Critical appraisal.} This is unusually close prior art: dynamics, uncertainty, balance, stopping and nonreactive crowds are already integrated. Its simplified walking and perception models differ from a learned G1 controller with measured latency, but those differences must become a defined new modeling or decision contribution. Passive safety also does not prevent a person from subsequently colliding with a stopped robot.

\paragraph{Implication for novel work.} This paper is a mandatory comparator and sharply narrows novelty. The defensible question is whether empirically calibrated policy response and acquisition-age uncertainty improve that safety-liveness tradeoff under actual onboard sensing, rather than whether those concepts should be combined at all.

\subsubsection{Barrier-function guarantees are conditional statements}\label{paper:cbf}
\textbf{Publication:} \citet{cbf}. \emph{Control Barrier Function Based Quadratic Programs for Safety Critical Systems}. \href{https://doi.org/10.1109/tac.2016.2638961}{Publication record}.

\textbf{Evidence anchor:} Selected full-text inspection. Author manuscript: abstract; Sections II-IV definitions and invariance discussion; automotive application framing. \href{https://arxiv.org/pdf/1609.06408}{Inspected source copy}.

\paragraph{Method and contribution.} Ames and colleagues formalize safety through forward invariance of a set and performance through Lyapunov-style objectives. Reciprocal and zeroing barrier constructions lead to admissible control conditions that can be enforced by a quadratic program. The optimization framework exposes conflicts between performance and safety instead of hiding both inside one undifferentiated reward. The mathematical guarantee belongs to the modeled system and stated regularity/feasibility conditions.

\paragraph{What the evidence establishes.} The paper develops theoretical results and illustrates the framework with adaptive-cruise-control and lane-keeping problems. Its proof structure is different from an empirical collision-rate comparison. Neither automotive examples nor an inequality written in a learning objective establish a humanoid implementation's model accuracy, sensing coverage or sampled-data behavior.

\paragraph{Critical appraisal.} A practical filter must address control limits, model error, observation uncertainty and whether a feasible safe action exists. A planner can produce a command that satisfies a nominal constraint and still fail if a later slew limiter or policy changes its realized trajectory. Replacing an online constraint with a finite-weight learning cost changes the assurance level and requires a separate argument.

\paragraph{Implication for novel work.} Use barrier literature to define precise guarantee language and to design a model-matched shield comparator. Geoff-Sim's current rollout supervisor should be described by its measured rejection behavior and assumptions unless a valid invariance argument is actually established.

\subsubsection{Learning robustness versus certifying safety}\label{paper:safecontrol}
\textbf{Publication:} \citet{safecontrol}. \emph{Safe Learning in Robotics: From Learning-Based Control to Safe Reinforcement Learning}. \href{https://doi.org/10.1146/annurev-control-042920-020211}{Publication record}.

\textbf{Evidence anchor:} Selected full-text inspection. Author survey: introduction and Section 3 taxonomy of learned dynamics, safe RL and certification. \href{https://arxiv.org/pdf/2108.06266}{Inspected source copy}.

\paragraph{Method and contribution.} Brunke and colleagues organize safe learning around prior model knowledge, uncertainty and the type of guarantee sought. The inspected taxonomy separates learning uncertain dynamics within a control framework, encouraging safe behavior through reinforcement learning, and certifying or modifying a learned controller through constraints or backup mechanisms. This distinguishes a performance-improving learned model from a validated safety envelope.

\paragraph{What the evidence establishes.} The work synthesizes control and learning methods, their assumptions and benchmark practices. It explains why guarantees depend on the formalized uncertainty and operating conditions. It does not certify an arbitrary controller merely because that controller uses a safety-related cost, a Bayesian model or a backup stop.

\paragraph{Critical appraisal.} The taxonomy highlights an important reporting failure: robustness demonstrated on sampled disturbances is often described with the language of formal safety. Conversely, formal guarantees can become operationally weak when their assumptions exclude the failures that matter. There is also a difference between keeping learning exploration safe and keeping a fixed deployed policy safe after training.

\paragraph{Implication for novel work.} Position a response-calibrated G1 method at the assurance level its evidence supports. Report model residual coverage, operating conditions and failure modes, then evaluate whether a stronger certified formulation remains feasible at the available compute and sensing budget.

\subsection{Learning interaction and predicting motion}
\subsubsection{CADRL: amortizing interaction reasoning into a value function}\label{paper:cadrl}
\textbf{Publication:} \citet{cadrl}. \emph{Decentralized non-communicating multiagent collision avoidance with deep reinforcement learning}. \href{https://doi.org/10.1109/icra.2017.7989037}{Publication record}.

\textbf{Evidence anchor:} Selected full-text inspection. Author preprint: abstract, robot-centric parameterization, value-function action selection and training algorithm. \href{https://arxiv.org/pdf/1609.07845}{Inspected source copy}.

\paragraph{Method and contribution.} CADRL learns an estimate of progress-to-goal from the joint observable state of interacting agents. At runtime, candidate velocity actions are scored through short lookahead and the learned value function. A robot-centric coordinate transformation removes translation/rotation redundancy. Neighbor goals are hidden, and the implementation approximates their immediate evolution using filtered velocity. Multiagent use builds on a policy learned around simpler interactions rather than solving a full joint trajectory optimization at every step.

\paragraph{What the evidence establishes.} The paper presents simulation comparisons with ORCA and reports improved travel-time quality under its protocol. The inspected formulation makes the assumed reciprocity and state-transition approximation visible. Its main computational advantage is shifting expensive interaction learning offline, not eliminating uncertainty about other agents' intentions.

\paragraph{Critical appraisal.} The observation interface already contains agent positions, velocities and radii, so detection, association and localization are separate responsibilities. Pairwise aggregation can miss collective effects, motivating later attention and graph methods. A learned state value is also not a physical feasibility certificate: candidate velocities still need a model of execution on the target robot. The meaning of the baseline comparison changes if a shared safety supervisor overrides most selected actions.

\paragraph{Implication for novel work.} CADRL supplies an interpretable learned-state baseline. Compare it on matched tracks, action limits and supervision, and report any mismatch between its training crowd behavior and the nonreactive or independently moving actors used for evaluation.

\subsubsection{SARL: learning which people matter collectively}\label{paper:sarl}
\textbf{Publication:} \citet{sarl}. \emph{Crowd-Robot Interaction: Crowd-Aware Robot Navigation With Attention-Based Deep Reinforcement Learning}. \href{https://doi.org/10.1109/icra.2019.8794134}{Publication record}.

\textbf{Evidence anchor:} Selected full-text inspection. arXiv v2 full HTML and PDF: Sections II-IV, implementation details and Tables I-II. \href{https://arxiv.org/pdf/1809.08835}{Inspected source copy}.

\paragraph{Method and contribution.} SARL embeds robot-human interactions and pools them with learned attention; LM-SARL additionally represents human-human interactions with local maps. The weighted embedding gives a fixed-size crowd representation despite varying neighbor count. A value network selects among sampled velocity actions. The paper explicitly assumes holonomic motion and instantaneous achievement of the commanded velocity, which is a materially different interface from a delayed walking policy.

\paragraph{What the evidence establishes.} The evaluation separates invisible robots, to which ORCA-controlled humans do not react, from visible robots with mutual influence, using 500 test cases in each setting. LM-SARL reports success 1.00 in the displayed tables; ORCA's invisible-setting success is 0.43. Real-world Segway demonstrations are also reported. These values describe a specific planar simulator and protocol, not complete sensing-to-motor reliability on a humanoid.

\paragraph{Critical appraisal.} The paper's controlled ablations support attention and local interaction encoding within that task. They do not make attention weights causal explanations of human intent. Training with true next-state information and deployment with an approximate transition model require careful distinction, as emphasized by RGL's later comparison. The 0.2 m discomfort metric is a separation proxy rather than a human-rating instrument.

\paragraph{Implication for novel work.} SARL is a strong architectural baseline for crowd aggregation, but its instantaneous-action assumption must be exposed in G1 comparisons. A contribution should show what calibrated execution and observation age add beyond better crowd-state encoding.

\subsubsection{Relational graphs and multi-step lookahead}\label{paper:rgl}
\textbf{Publication:} \citet{rgl}. \emph{Relational Graph Learning for Crowd Navigation}. \href{https://doi.org/10.1109/iros45743.2020.9340705}{Publication record}.

\textbf{Evidence anchor:} Selected full-text inspection. Author preprint: model/training description, Table I, planning-budget discussion and physical demonstration description. \href{https://arxiv.org/pdf/1909.13165}{Inspected source copy}.

\paragraph{Method and contribution.} RGL learns interactions among agents through a graph convolutional representation shared by state prediction and value estimation. Predicting future states enables multi-step lookahead instead of relying solely on a one-step terminal value. The planning depth and retained action width create an explicit tradeoff between computation and anticipation. This is a learned-model planning approach, distinct from directly regressing the next command from sensor history.

\paragraph{What the evidence establishes.} Table I reports five independently seeded training runs and 500 test cases for a five-person circle crossing. MP-RGL-Multistep reports success 0.96 plus/minus 0.02 and collision 0.02 plus/minus 0.01 under that protocol. The comparison replaces a ground-truth-next-state assumption with a linear prediction variant before assessing the learned relation/model contribution. Pioneer demonstrations provide additional qualitative deployment context.

\paragraph{Critical appraisal.} The source-contract adjustment is central: architecture gains cannot be separated from privileged transition information otherwise. Additional lookahead can improve behavior but also magnify model errors and consume the time needed to react. An average planning time alone does not capture this feedback. The physical demonstrations do not provide a matched biped dynamics or latency study.

\paragraph{Implication for novel work.} Use RGL to motivate joint reporting of prediction fidelity, planning horizon, computation and resulting mission outcomes. A G1 comparison should hold the observation contract and search budget fixed before attributing gains to response calibration or graph structure.

\subsubsection{Temporal crowd reasoning under restricted visibility}\label{paper:dsrnn}
\textbf{Publication:} \citet{dsrnn}. \emph{Decentralized Structural-RNN for Robot Crowd Navigation with Deep Reinforcement Learning}. \href{https://doi.org/10.1109/icra48506.2021.9561595}{Publication record}.

\textbf{Evidence anchor:} Selected full-text inspection. Retrieved arXiv v4: experiment setup, field-of-view/group evaluation, real-world section and limitations. \href{https://arxiv.org/pdf/2011.04820}{Inspected source copy}.

\paragraph{Method and contribution.} DS-RNN represents spatial and temporal relationships through structured recurrent components and trains a model-free navigation policy. Memory helps the robot act when neighbors are not fully observable, while the structured graph assigns different roles to robot history and robot-human interactions. The method avoids requiring an explicit analytic pedestrian transition model or expert initialization as the central training mechanism.

\paragraph{What the evidence establishes.} The inspected evaluation varies field of view across 90, 180 and 360 degrees and group populations across 10, 15 and 20 humans, with 500 unseen cases per setting. A recurrent-attention ablation uses matched PPO settings to better isolate the architecture. Physical TurtleBot 2i demonstrations use RGB-D detection/tracking and a tracking camera, with processing on a remote host.

\paragraph{Critical appraisal.} Restricted field of view is not equivalent to arbitrary detector failure, stale tracks or calibrated uncertainty. The limitations explicitly recognize the difference between an invisible simulated robot and reciprocal real interactions. Remote processing also changes the deployment timing boundary. The retrieved preprint is a later version than the conference year; its version must be retained in any detailed comparison.

\paragraph{Implication for novel work.} A memory-based comparator prevents the proposed project from treating partial observability as an untouched problem. Test controlled latency and track-age interventions separately from visibility, and distinguish temporal representation gains from low-level response-model gains.

\subsubsection{A taxonomy of what prediction models actually assume}\label{paper:prediction}
\textbf{Publication:} \citet{prediction}. \emph{Human motion trajectory prediction: a survey}. \href{https://doi.org/10.1177/0278364920917446}{Publication record}.

\textbf{Evidence anchor:} Selected full-text inspection. Author survey: taxonomy/classification rules and physics-based, pattern-based and contextual-model discussion. \href{https://arxiv.org/pdf/1905.06113}{Inspected source copy}.

\paragraph{Method and contribution.} Rudenko and colleagues organize trajectory prediction around motion models, contextual information and other modeling choices. Physics-based methods encode explicit dynamics; pattern-based methods learn regularities from observed motion; higher-level reasoning can incorporate goals, semantics and interactions. Context is not a single binary attribute: obstacle awareness, map awareness and semantic awareness provide different information and failure modes.

\paragraph{What the evidence establishes.} The survey synthesizes methods and datasets across robotics and related communities. Its classification rules explicitly acknowledge mixed methods rather than forcing every approach into one exclusive category. The result is a vocabulary for choosing a forecast appropriate to the task and data, not evidence that the most complex predictor is universally preferable.

\paragraph{Critical appraisal.} A forecast is evaluated under a conditioning set. A model supplied with a person's destination cannot be ranked directly against a model using only recent detections. Long-horizon trajectory error, immediate collision coverage and response to a robot intervention are different targets. Survey breadth also means that detailed claims about an individual predictor should be checked against its primary paper.

\paragraph{Implication for novel work.} Specify prediction horizon, sensor history, map/goal information and whether the robot's planned action conditions the prediction. Constant velocity should remain an explicit low-cost baseline, with failures analyzed by maneuver and observation quality rather than dismissed categorically.

\subsubsection{Social LSTM: recurrent interaction encoding before planning}\label{paper:sociallstm}
\textbf{Publication:} \citet{sociallstm}. \emph{Social LSTM: Human Trajectory Prediction in Crowded Spaces}. \href{https://doi.org/10.1109/cvpr.2016.110}{Publication record}.

\textbf{Evidence anchor:} Selected full-text inspection. CVPR paper: problem description, social-pooling formulation, Section 4 and Table 1. \href{https://openaccess.thecvf.com/content_cvpr_2016/papers/Alahi_Social_LSTM_Human_CVPR_2016_paper.pdf}{Inspected source copy}.

\paragraph{Method and contribution.} Social LSTM models each pedestrian's history with a recurrent network and pools neighboring hidden states to represent local interaction. It predicts future pedestrian motion from tracked positions, allowing curved and mutually influenced trajectories rather than independent straight-line extrapolation. The social pooling operation is a learned representation of interaction, not a navigation policy or a guarantee that a robot action will be acceptable.

\paragraph{What the evidence establishes.} The paper evaluates five ETH/UCY subsets and reports displacement metrics over a 4.8-second forecast. Table 1 notes that the IGP comparator receives the intended destination, unlike other methods. Improvements are not uniform: the average final-displacement entry for Social LSTM is 0.61 versus 0.60 for the listed social-force baseline, despite advantages on other metrics and scenes.

\paragraph{Critical appraisal.} The paper illustrates why an abstract-level claim of outperforming prior work needs table-level qualification. Dataset splits, synthetic pretraining, definitions of displacement error and privileged destinations affect comparison. Predicting trajectories from already available tracks does not solve missing detections or identity swaps. A robot also needs to turn the prediction into a feasible decision before its scene changes.

\paragraph{Implication for novel work.} Use the method to establish interaction-aware prediction history, not to claim solved closed-loop navigation. Evaluate forecast-to-planner coupling and collision-critical errors, and preserve the conditioning information and metric definition in any reproduced baseline.

\subsubsection{Multimodal futures and the limits of best-of-sample evaluation}\label{paper:socialgan}
\textbf{Publication:} \citet{socialgan}. \emph{Social GAN: Socially Acceptable Trajectories with Generative Adversarial Networks}. \href{https://doi.org/10.1109/cvpr.2018.00240}{Publication record}.

\textbf{Evidence anchor:} Selected full-text inspection. Author preprint: model overview, Section 4 protocol, sampling/variety-loss analysis and pooling discussion. \href{https://arxiv.org/pdf/1803.10892}{Inspected source copy}.

\paragraph{Method and contribution.} Social GAN uses a recurrent encoder-decoder generator, interaction pooling and an adversarial discriminator to produce several plausible futures from the same observed history. A variety loss encourages at least one sample to match the recorded future. This addresses the inadequacy of averaging mutually exclusive routes into one implausible mean trajectory, while retaining efficient joint prediction across pedestrians.

\paragraph{What the evidence establishes.} ETH/UCY evaluation uses leave-one-out scenes, eight observed steps and forecasts of eight or twelve steps at 0.4-second intervals. Quantitative reporting selects the best sampled prediction in an L2 sense. The paper also reports difficulty reproducing the original Social LSTM ranking without its synthetic-data setup, and notes that pooling can improve qualitative coherence without improving every displacement metric.

\paragraph{Critical appraisal.} Best-of-N error benefits from increasing sample count and does not reveal which sample a robot should trust before seeing the future. Diversity is not a calibrated probability distribution over dangerous outcomes. A planner that rejects every sampled path can freeze, while selecting the convenient sample can understate risk. The timing cost of generating and checking samples belongs in the closed-loop evaluation.

\paragraph{Implication for novel work.} Test whether multimodality improves encounter coverage and collision-free progress under a matched computational budget. Compare against constant velocity and calibrated occupancy regions rather than equating a lower best-of-N score with safer navigation.

\subsubsection{Trajectron++: dynamics-aware probabilistic forecasting}\label{paper:trajectron}
\textbf{Publication:} \citet{trajectron}. \emph{Trajectron++: Dynamically-Feasible Trajectory Forecasting with Heterogeneous Data}. \href{https://doi.org/10.1007/978-3-030-58523-5_40}{Publication record}.

\textbf{Evidence anchor:} Selected full-text inspection. Author preprint: desiderata and architecture description; Section 5 datasets, metrics and comparator definitions. \href{https://arxiv.org/pdf/2001.03093}{Inspected source copy}.

\paragraph{Method and contribution.} Trajectron++ combines a graph-structured recurrent representation with latent-variable forecasting, heterogeneous inputs such as semantic maps, and agent-specific dynamics. Its predictions can be conditioned on a proposed ego trajectory, making the robot's action part of the forecasting problem. This is a more appropriate conceptual interface for interaction-aware planning than assuming that each person's future is independent of what the robot does.

\paragraph{What the evidence establishes.} Evaluation spans pedestrian ETH/UCY data and heterogeneous nuScenes traffic. The inspected protocol reports average/final displacement, KDE-based negative log likelihood and best-of-N metrics. Those measures assess different properties of a generative model. The examples demonstrate a flexible forecasting framework, not a closed-loop G1 contact/comfort study.

\paragraph{Critical appraisal.} Dynamics-consistent forecasts are not necessarily probability-calibrated under a new domain. Ego-conditioned prediction also depends on the distribution of ego actions represented in the training data: an implausible proposed maneuver can induce an unreliable response forecast. Map availability, object classes and sampling frequency differ considerably from a torso-mounted LiDAR tracker in an online-built map.

\paragraph{Implication for novel work.} The paper rules out novelty based simply on combining forecasting with dynamics or ego-plan conditioning. The open question is whether uncertainty and sensing age can be connected to executable humanoid decisions with verified coverage and useful progress.

\subsection{Perception and state estimation}
\subsubsection{Learning range-data detections instead of hand-crafting clusters}\label{paper:drow}
\textbf{Publication:} \citet{drow}. \emph{DROW: Real-Time Deep Learning-Based Wheelchair Detection in 2-D Range Data}. \href{https://doi.org/10.1109/lra.2016.2645131}{Publication record}.

\textbf{Evidence anchor:} Selected full-text inspection. Author preprint: Section II preprocessing/voting and Section III dataset and split design. \href{https://arxiv.org/pdf/1603.02636}{Inspected source copy}.

\paragraph{Method and contribution.} DROW builds distance-normalized cutouts around individual laser returns, predicts a class and an offset toward an object's center, and consolidates these votes using a grid and non-maximum suppression. This separates the sensor surface points from the latent object center. It also addresses the changing angular extent of an object with range, which makes a fixed-width raw scan window problematic.

\paragraph{What the evidence establishes.} The original paper focuses on wheelchairs and walkers, with approximately 464,000 scans and 24,012 annotated scans from a care facility. Training, validation and test data include spatial separation and temporally disjoint sequences. The described scanner is mounted near the floor, around 37 cm, a key detail for assessing transfer to humanoid sensing.

\paragraph{Critical appraisal.} The method is not originally a universal person detector, even though later DROW work extends that scope. Dense scan counts should not be mistaken for equally many independent environments or people. A detector trained on low-mounted 2D geometry may not tolerate torso-height 3D projections, partial upper-body returns or walking-induced sensor motion without adaptation.

\paragraph{Implication for novel work.} Compare learned center estimation with geometric clustering using matched projections and ground truth. Measure bias, misses and false positives separately, and verify that perception changes improve downstream navigation rather than only a detector score.

\subsubsection{Temporal range features with bounded per-frame work}\label{paper:drspaam}
\textbf{Publication:} \citet{drspaam}. \emph{DR-SPAAM: A Spatial-Attention and Auto-regressive Model for Person Detection in 2D Range Data}. \href{https://doi.org/10.1109/iros45743.2020.9341689}{Publication record}.

\textbf{Evidence anchor:} Selected full-text inspection. Author preprint: Section III temporal aggregation; Section IV, Table I and dataset split description. \href{https://arxiv.org/pdf/2004.14079}{Inspected source copy}.

\paragraph{Method and contribution.} DR-SPAAM retains a feature template and updates it as new scans arrive. Spatial attention associates nearby features across time, while an autoregressive update accumulates evidence without repeatedly rebuilding cutouts for every previous scan. Its forward-looking aggregation contrasts with backward alignment/recomputation and aims to preserve the benefit of temporal observations at lower computational cost.

\paragraph{What the evidence establishes.} The evaluation uses the DROW dataset's person annotations and explicitly separates retrained baselines from previously reported DROW variants. Table I reports AP at a 0.5 m association threshold of 70.3 for DR-SPAAM and 67.9 for the retrained five-scan baseline. Ablations remove spatial attention or autoregression, giving a more informative test than comparing only complete systems.

\paragraph{Critical appraisal.} The attribution depends on using the correct retrained baseline and association criterion. Faster network inference is not automatically lower observation-to-command age after ROS transport, queues and planning. A recurrent template can also retain misleading features; deployment needs a policy for missing scans, reset and stale evidence. The sensor-domain difference from humanoid LiDAR remains important.

\paragraph{Implication for novel work.} The architecture offers a concrete lesson for efficient temporal perception. Evaluate throughput, retained-feature age, center/velocity error and navigation outcomes together, using the actual sensor mounting and available compute.

\subsubsection{Direct LiDAR-inertial estimation and incremental mapping}\label{paper:fastlio}
\textbf{Publication:} \citet{fastlio}. \emph{FAST-LIO2: Fast Direct LiDAR-Inertial Odometry}. \href{https://doi.org/10.1109/tro.2022.3141876}{Publication record}.

\textbf{Evidence anchor:} Selected full-text inspection. Author preprint: contribution/estimation discussion and Section VI data-structure benchmark setup. \href{https://arxiv.org/pdf/2107.06829}{Inspected source copy}.

\paragraph{Method and contribution.} FAST-LIO2 couples inertial propagation with iterated filtering and direct registration of raw LiDAR returns to a map, avoiding a mandatory hand-engineered feature extractor. An incremental k-d tree supports insertion, deletion, rebalancing and on-tree downsampling. The combination addresses both estimation and the data-structure cost of maintaining useful geometry at high update rates.

\paragraph{What the evidence establishes.} The paper separates data-structure evaluation from whole-system odometry comparisons and tests varied LiDAR/IMU datasets and real motions. The inspected setup explains how alternative map structures are integrated and downsampled for comparison. The abstract's fast-update claims therefore belong to identified datasets and hardware, not to an arbitrary ROS graph containing the estimator.

\paragraph{Critical appraisal.} This is a state-estimation contribution, not a social navigation controller. Reliable deskew and fusion require acquisition timing, extrinsics and an appropriate IMU contract. A simulator that raycasts a frozen instantaneous scene does not automatically test rolling-scan distortion. Improved trajectory error also does not by itself solve stale person tracks or command-response mismatch.

\paragraph{Implication for novel work.} Treat FAST-LIO2 as a controlled estimator alternative. Match sensor sampling and calibration, then compare both localization metrics and changes in navigation decisions. Keep robot-pose truth isolated from onboard control in every experimental arm.

\subsubsection{Visual-inertial initialization and persistent map reuse}\label{paper:orbslam}
\textbf{Publication:} \citet{orbslam}. \emph{ORB-SLAM3: An Accurate Open-Source Library for Visual, Visual–Inertial, and Multimap SLAM}. \href{https://doi.org/10.1109/tro.2021.3075644}{Publication record}.

\textbf{Evidence anchor:} Selected full-text inspection. Accepted author manuscript: overview of visual/inertial/multi-map contributions and MAP/IMU-initialization formulation. \href{https://arxiv.org/pdf/2007.11898}{Inspected source copy}.

\paragraph{Method and contribution.} ORB-SLAM3 combines feature-based visual estimation, tightly integrated inertial measurements and multiple-map management. The inspected formulation combines visual reprojection and inertial residuals, with uncertainty-aware maximum-a-posteriori initialization of scale, gravity, velocity and biases. Place recognition and map reuse extend the objective beyond short-term visual odometry.

\paragraph{What the evidence establishes.} The manuscript supports several camera configurations and explicitly discusses initialization and data association. Those technical details explain why importing an estimator means more than subscribing to a pose topic. This appraisal does not restate a cross-dataset accuracy multiplier as though it were a measured improvement for Geoff-Sim.

\paragraph{Critical appraisal.} Visual observability, calibration, exposure and motion characteristics differ from the project's LiDAR-centric pipeline. Map reuse can help recovery but introduces coordinate corrections that downstream tracks, goals and commands must handle consistently. Dataset trajectory accuracy is not an end-to-end measure of crowd avoidance, and a camera-rich estimator cannot be fairly compared without accounting for its different sensors.

\paragraph{Implication for novel work.} Use ORB-SLAM3 to frame alternative sensing and map-lifecycle experiments, not as a presumed drop-in upgrade. The scientific question is which estimation errors actually change collision-critical decisions under matched information and latency budgets.

\subsection{Datasets, benchmarks and social validity}
\subsubsection{Egocentric perception exposes a different failure distribution}\label{paper:jrdb}
\textbf{Publication:} \citet{jrdb}. \emph{JRDB: A Dataset and Benchmark of Egocentric Robot Visual Perception of Humans in Built Environments}. \href{https://doi.org/10.1109/tpami.2021.3070543}{Publication record}.

\textbf{Evidence anchor:} Selected full-text inspection. Author preprint: dataset overview and detection/tracking results, including Tables 3-4 and occlusion discussion. \href{https://arxiv.org/pdf/1910.11792}{Inspected source copy}.

\paragraph{Method and contribution.} JRDB records multimodal observations from a robot's viewpoint and provides human detection/tracking annotations. Panoramic imagery, LiDAR and associated trajectories expose occlusion, moving sensors and dense pedestrian configurations that differ from conventional image collections. Its evaluation decomposes detection quality and tracking consistency instead of assuming that a planner receives an accurate set of human states.

\paragraph{What the evidence establishes.} The inspected results use AP and OSPA-style errors for detection, with tracking measures that include false positives, false negatives and identity changes. The discussion shows that methods can rank differently on localization and cardinality error. It also analyzes adaptation to panoramic images, training on JRDB versus generic pretraining, and the sensitivity of results to occlusion.

\paragraph{Critical appraisal.} Perception metrics have to be interpreted jointly: a detector that localizes its surviving detections well can still miss many people. Large annotation counts do not eliminate environment or sensor-domain bias. The dataset does not itself establish how a G1's controller responds to those errors, and its sensor arrangement differs from a walking torso-mounted scanner.

\paragraph{Implication for novel work.} Borrow the decomposition of missing, misplaced and misassociated people. A navigation evaluation should connect those errors to encounter timing and applied commands, rather than reporting only nearest-track distance or treating detection count as recall.

\subsubsection{Demonstration data and a genuine human-evaluation precedent}\label{paper:scand}
\textbf{Publication:} \citet{scand}. \emph{Socially CompliAnt Navigation Dataset (SCAND): A Large-Scale Dataset of Demonstrations for Social Navigation}. \href{https://doi.org/10.1109/lra.2022.3184025}{Publication record}.

\textbf{Evidence anchor:} Selected full-text inspection. Author manuscript: dataset description, local-planner evaluation summary/Figure 7, anticipated use cases and conclusion. \href{https://arxiv.org/pdf/2203.15041}{Inspected source copy}.

\paragraph{Method and contribution.} SCAND captures human-teleoperated navigation with multimodal sensors and commands on morphologically different robots, including Spot and Jackal. Demonstrations supply behavioral examples when a reward function is hard to specify. The paper also explores demonstrator classification and imitation of global/local planning behavior, recognizing that multiple human strategies can be socially plausible.

\paragraph{What the evidence establishes.} The recorded dataset totals 8.7 hours, 138 trajectories and 25 miles. The inspected evaluation reports a study with fourteen human participants for a learned local planner in two scenarios, alongside the demonstration data. It therefore directly contradicts a blanket claim that social-navigation research has no human-subject evaluation.

\paragraph{Critical appraisal.} The authors identify geographic and convention bias from collection in Austin, such as passing preferences. Four demonstrators and a finite set of scenarios cannot represent all users or norms. The local-planner study is more relevant to perceived behavior than a pure clearance metric, but it does not establish universal comfort or transfer to a humanoid carrying a payload.

\paragraph{Implication for novel work.} Use SCAND as both an imitation-data reference and a methodological comparator for human judgments. The narrower open question is how humanoid morphology, motion transients and task context affect perceived comfort under a controlled navigation comparison.

\subsubsection{Grounded scenarios, common metrics and baseline adaptation}\label{paper:socnavbench}
\textbf{Publication:} \citet{socnavbench}. \emph{SocNavBench: A Grounded Simulation Testing Framework for Evaluating Social Navigation}. \href{https://doi.org/10.1145/3476413}{Publication record}.

\textbf{Evidence anchor:} Selected full-text inspection. Author preprint: framework overview, Section 4 algorithms/protocol and Section 5 metric interpretation. \href{https://arxiv.org/pdf/2103.00047}{Inspected source copy}.

\paragraph{Method and contribution.} SocNavBench combines curated scenarios grounded in recorded pedestrian data with a shared simulation and evaluation interface. It supports comparing navigation algorithms through several outcome dimensions instead of a single success scalar. Realistic rendering is a capability of the platform, while the experimental modes and information supplied to an algorithm remain separately configurable.

\paragraph{What the evidence establishes.} The inspected comparison includes Social Forces, ORCA, SA-CADRL and a pedestrian-unaware planner. Because pedestrians do not react to the robot, the authors adjust ORCA collaboration coefficients. They explicitly discuss kinematic constraints, the use of a meta-planner and selecting only the nearest pedestrians for a fixed-input network. These implementation decisions materially affect what the benchmark tests.

\paragraph{Critical appraisal.} A common environment does not automatically mean an equal information or dynamics contract. Recorded pedestrian replay gives repeatability but removes counterfactual reactions to changed robot behavior. Synchronous evaluation can also hide latency effects unless timing is modeled separately. The paper's multidimensional metric analysis is stronger than extracting one aggregate ranking while ignoring these adaptations.

\paragraph{Implication for novel work.} Geoff-Sim's benchmark contribution must add a new controlled question, such as delay/response provenance or articulated contact, beyond merely providing worlds and a GUI. Document baseline adapters, actor responsiveness and timing modes with the result.

\subsubsection{Configurable multi-robot learning in constrained shared spaces}\label{paper:socialgym}
\textbf{Publication:} \citet{socialgym}. \emph{SOCIALGYM 2.0: Simulator for Multi-Robot Learning and Navigation in Shared Human Spaces}. \href{https://doi.org/10.1609/aaai.v38i21.30562}{Publication record}.

\textbf{Evidence anchor:} Selected full-text inspection. AAAI paper: simulator architecture, observations/rewards, environment configuration and limitations. \href{https://ojs.aaai.org/index.php/AAAI/article/download/30562/32724}{Inspected source copy}.

\paragraph{Method and contribution.} SocialGym 2.0 emphasizes multiagent learning and configurable indoor interactions such as doorways, corridors, intersections and roundabouts. Observation and reward functions can be customized, along with environmental dimensions, density and robot kinodynamic settings. Its configuration-oriented design makes experimental variation a first-class feature instead of tying every method to one hard-coded social objective.

\paragraph{What the evidence establishes.} The short system paper describes extensibility and integration with multiagent learning interfaces. It explicitly notes CPU/parallel-processing limitations and lists continuous actions and further MARL algorithms as future development. Those statements constrain what may be inferred from the platform's advertised flexibility.

\paragraph{Critical appraisal.} Configurability is valuable, but each configured benchmark still needs validation. Changing rewards, state observations and dynamics together can undermine algorithm comparisons. Multi-robot interaction is not automatically a model of uncooperative human motion. A simulator architecture paper also supplies a different type of evidence from a statistically characterized physical deployment.

\paragraph{Implication for novel work.} Adopt explicit scenario/configuration variation and retain its provenance. Position Geoff-Sim's distinct contribution around measured biped execution and source-audited evaluation, rather than claiming that configurable shared-space benchmarks or mixed dynamics are absent from existing tools.

\subsubsection{Compositional urban variation and the success-cost tradeoff}\label{paper:metaurban}
\textbf{Publication:} \citet{metaurban}. \emph{MetaUrban: An Embodied AI Simulation Platform for Urban Micromobility}. \href{https://arxiv.org/abs/2407.08725}{Publication record}.

\textbf{Evidence anchor:} Selected full-text inspection. arXiv v2: platform overview, Section 4 tasks/protocol and Table 1 benchmarks. \href{https://arxiv.org/pdf/2407.08725}{Inspected source copy}.

\paragraph{Method and contribution.} MetaUrban procedurally composes urban spaces, obstacles, moving people and heterogeneous mobility platforms. Its benchmark separates static point navigation from social navigation and compares online RL, safe RL, offline RL and imitation approaches. The inspected main experiment uses a wheeled delivery platform with acceleration, braking and steering actions, rather than a low-level G1 locomotion policy.

\paragraph{What the evidence establishes.} Table 1 reports in-distribution and unseen-environment outcomes, including success, path/social metrics and cumulative costs. For the displayed social-navigation test, PPO has higher success than its constrained variants while those variants can reduce cost. The report explicitly discusses the difference between improving safety properties and maintaining effectiveness.

\paragraph{Critical appraisal.} The large scene library and asset count do not themselves establish behavioral realism or coverage of rare interaction failures. Generalization to unseen layouts is not identical to generalization to a new controller, sensor mounting or crowd policy. Cost reductions should be interpreted together with mission completion and episode exposure, especially when a method stops or fails early.

\paragraph{Implication for novel work.} Use compositional scene variation and held-out layouts, but retain the lower-level G1 execution question. Report a collision-completion frontier rather than treating either safety cost or success alone as the definitive result.

\subsubsection{Human-in-the-loop tasks and the meaning of simulated embodiment}\label{paper:habitat3}
\textbf{Publication:} \citet{habitat3}. \emph{Habitat 3.0: A Co-Habitat for Humans, Avatars and Robots}. \href{https://arxiv.org/abs/2310.13724}{Publication record}.

\textbf{Evidence anchor:} Selected full-text inspection. Author report: Section 3 avatar/simulator design and social-navigation baseline/sensor-ablation discussion. \href{https://arxiv.org/pdf/2310.13724}{Inspected source copy}.

\paragraph{Method and contribution.} Habitat 3.0 combines articulated human avatars, efficient appearance/motion generation and interfaces for real people to interact with simulated robots. Social navigation includes finding and following a person, while social rearrangement introduces collaborative task structure. The paper separates skeletal collision representation from skinned visual appearance, which is important when interpreting the physical meaning of a rendered scene.

\paragraph{What the evidence establishes.} The inspected navigation experiment uses Spot in train/validation/test scene splits and compares a map-informed heuristic expert with sensor-to-action RL and sensor ablations. The stated RL inputs include depth, a humanoid detector and humanoid GPS. Consequently, the observation contract must be read before calling a method fully unprivileged simply because it also consumes images.

\paragraph{Critical appraisal.} Human-in-the-loop simulation can capture responses that prerecorded replay cannot, but it is not identical to in-person interaction with a physical robot. Avatar articulation also does not imply that its gait is a dynamically balanced low-level humanoid controller. Task definitions such as person following differ from independent waypoint avoidance and need distinct success criteria.

\paragraph{Implication for novel work.} Habitat provides strong prior art for interactive evaluation and collaborative tasks. Geoff-Sim should claim its specific physics/control and provenance contribution, and use human-in-the-loop tools as a possible evaluation bridge rather than inferring hardware validity from visual realism.

\subsubsection{Composed tasks reveal the gap between skills and missions}\label{paper:omninav}
\textbf{Publication:} \citet{omninav}. \emph{Beyond Isolation: A Unified Benchmark for General-Purpose Navigation}. \href{https://arxiv.org/abs/2605.09441v1}{Publication record}.

\textbf{Evidence anchor:} Selected full-text inspection. arXiv v1: evaluation platform, Tables II-IV and task/metric definitions. \href{https://arxiv.org/pdf/2605.09441v1}{Inspected source copy}.

\paragraph{Method and contribution.} OmniNavBench composes navigation and interaction objectives across different embodiments in an Isaac Sim platform with continuous physical stepping and locomotion controllers. Its tasks include goal-directed navigation, following and question answering, with social constraints applied across tasks. This extends evaluation from an isolated skill to sequential mission execution.

\paragraph{What the evidence establishes.} The inspected tables distinguish whole-episode success, subtask completion, continuous progress, following and social intrusion. They use different completion predicates by task; a PointNav subtask's distance tolerance and the final-episode stopping criterion are not interchangeable. The published comparison includes Carter, Aliengo and H1, and shows that partial progress can coexist with very low complete-task success.

\paragraph{Critical appraisal.} A personal-space violation ratio is not a count of articulated physical contacts. The platform's reproducibility claims also require reproducing actual initialization, actor schedules and execution, not just retaining a seed. Composite task failures may arise from language, perception, navigation or task transitions; a low total success rate does not identify the bottleneck.

\paragraph{Implication for novel work.} The project cannot claim the first physically grounded multi-task navigation benchmark. A meaningful extension would diagnose why complete humanoid missions fail through independent source, timing, control and lifecycle instrumentation, with task-specific validity and stopping criteria.

\subsubsection{Dynamic-pedestrian language navigation and metric conflicts}\label{paper:dped}
\textbf{Publication:} \citet{dped}. \emph{DPed-VLN: A Benchmark for Socially Compliant Vision-and-Language Navigation in Dynamic Pedestrian Environments}. \href{https://arxiv.org/abs/2609.21504}{Publication record}.

\textbf{Evidence anchor:} Selected full-text inspection. September 2026 arXiv v1: dataset and protocol overview, Tables 2-4 and pedestrian-dynamics ablation. \href{https://arxiv.org/pdf/2609.21504}{Inspected source copy}.

\paragraph{Method and contribution.} DPed-VLN couples navigation episodes with ordinary route instructions and prior-augmented language about dynamic pedestrians. ORCA-controlled humanoid actors populate Habitat 3.0 scenes, and the evaluation separates goal completion, efficiency and human-related metrics. The language-conditioned problem is adjacent to, but different from, a LiDAR/Nav2 waypoint controller.

\paragraph{What the evidence establishes.} The paper reports 33,093 episodes and paired instruction conditions. Its tables compare imitation and reinforcement learning and show that stronger success does not imply a uniformly lower human-collision rate. The static-versus-dynamic pedestrian ablation exposes a major failure distribution shift even when goal-success changes appear modest.

\paragraph{Critical appraisal.} The benchmark is a recent preprint, and the reported baselines/protocol are not independently replicated here. Humanoid pedestrians should not be confused with a low-level humanoid robot policy. Prior-augmented instructions add information that an ordinary goal-only system does not receive. Collision and personal-space metrics need their exact event definitions before comparison with a physics-contact scorer.

\paragraph{Implication for novel work.} This is evidence that dynamic human environments and broader task grounding are active research areas, not untouched novelty territory. The useful lesson for Geoff-Sim is to separate source privilege, task success and social-risk metrics and to test dynamic exposure explicitly.

\subsection{Locomotion, embodiment and transfer}
\subsubsection{Rapid adaptation changes the locomotion interface}\label{paper:rma}
\textbf{Publication:} \citet{rma}. \emph{RMA: Rapid Motor Adaptation for Legged Robots}. \href{https://doi.org/10.15607/rss.2021.xvii.011}{Publication record}.

\textbf{Evidence anchor:} Selected full-text inspection. Author preprint: overview of base policy/adaptation, simulation and hardware evaluation, adaptation analysis and conclusion. \href{https://arxiv.org/pdf/2107.04034}{Inspected source copy}.

\paragraph{Method and contribution.} RMA combines a locomotion policy with an adaptation module that infers relevant extrinsic conditions from recent experience. Training can use privileged environment parameters, while deployment must infer their effect from onboard histories. The architecture aims to adapt rapidly to changing payload, friction and terrain rather than requiring a new policy for every condition. This makes locomotion response a state-dependent property rather than a fixed velocity gain.

\paragraph{What the evidence establishes.} The paper separates simulated parameter variation from physical payload, step, compliant-surface and outdoor tests. Its simulation table averages three initialized policies and 1,000 episodes per initialization. Hardware trials are much smaller, with some severely failing methods stopped after two attempts to limit damage. Removing adaptation degrades the reported performance.

\paragraph{Critical appraisal.} The hardware protocol and simulation denominator cannot be pooled. Staying upright without progressing is not equivalent to successful locomotion, a distinction visible in the ablation discussion. RMA is primarily proprioceptive and does not by itself plan around people. Adaptation can also change the response model expected by a higher-level planner, so a static calibration may become stale.

\paragraph{Implication for novel work.} Evaluate whether a navigation response model must condition on recent policy state, load or terrain. Compare calibrated fixed-policy behavior with adaptation explicitly, rather than treating the low-level policy as an invariant actuator.

\subsubsection{Massive parallelism is an experimental enabler, not an acceptance result}\label{paper:rudin}
\textbf{Publication:} \citet{rudin}. \emph{Learning to Walk in Minutes Using Massively Parallel Deep Reinforcement Learning}. \href{https://arxiv.org/abs/2109.11978}{Publication record}.

\textbf{Evidence anchor:} Selected full-text inspection. Retrieved author preprint v3: training overview, Section 4 parallelism experiments, simulation setup and conclusion. \href{https://arxiv.org/pdf/2109.11978}{Inspected source copy}.

\paragraph{Method and contribution.} Rudin and colleagues combine an end-to-end GPU simulation/training pipeline with a terrain curriculum and training choices appropriate to thousands of parallel robots. The work analyzes how environment count, per-environment rollout length and batch size interact. Its contribution is not simply running more copies: insufficient temporal horizon can degrade learning even when total sample count is held fixed.

\paragraph{What the evidence establishes.} The study varies parallel robot counts and batch sizes and reports multiple training runs. It identifies a tradeoff region rather than unlimited monotonic benefit from parallelism. The conclusion explicitly states that the aim is fast generation of useful locomotion policies, not the absolute highest robustness achievable with every additional technique.

\paragraph{Critical appraisal.} Wall-clock training time is conditional on GPU, simulator, task and selected performance threshold. A quickly trained policy can be a valuable research baseline without establishing rare-failure reliability or social competence. Increasing the number of correlated simulator instances also does not create diverse independent deployment environments by itself.

\paragraph{Implication for novel work.} Use efficient training to support controlled ablations, repetitions and explicit initialization-only baselines. For payload work, distinguish gains from architecture, initialization, training duration and integration corrections before attributing success to reinforcement learning.

\subsubsection{Belief-state locomotion when exteroception is unreliable}\label{paper:rough}
\textbf{Publication:} \citet{rough}. \emph{Learning robust perceptive locomotion for quadrupedal robots in the wild}. \href{https://doi.org/10.1126/scirobotics.abk2822}{Publication record}.

\textbf{Evidence anchor:} Selected full-text inspection. Author manuscript: recurrent-belief explanation, field deployment results, Figure 4 and discussion/method overview. \href{https://arxiv.org/pdf/2201.08117}{Inspected source copy}.

\paragraph{Method and contribution.} Miki and colleagues integrate proprioceptive and exteroceptive observations through a recurrent belief representation. A privileged teacher supports training a deployable controller that exploits terrain information when it is reliable and falls back toward proprioceptive competence when observations are missing or misleading. This addresses the practical problem that visual geometry and physical support are not equivalent, for example with snow, vegetation or transparent obstacles.

\paragraph{What the evidence establishes.} The paper reports varied outdoor and underground deployments and inspects internal belief estimates under sensory failure. Examples include covering sensors, encountering compliant support and revising terrain estimates after contact. These experiments are more informative about sensing robustness than a clean height-map benchmark alone.

\paragraph{Critical appraisal.} Contact can legitimately disambiguate whether terrain supports a foot; using contact to resolve uncertainty about a nearby person has a different consequence. The task is terrain locomotion rather than human-centered navigation, and the amount of deployment exposure matters when interpreting zero-fall statements. A recurrent representation also does not automatically expose a calibrated uncertainty estimate to a separate planner.

\paragraph{Implication for novel work.} Use selective trust in perception as a design principle, but test the human-avoidance implications separately. The project needs an explicit observation-age and uncertainty interface, not merely a hidden state whose apparent robustness is assumed to transfer.

\subsubsection{Planning with a repertoire of executable skills}\label{paper:parkour}
\textbf{Publication:} \citet{parkour}. \emph{ANYmal parkour: Learning agile navigation for quadrupedal robots}. \href{https://doi.org/10.1126/scirobotics.adi7566}{Publication record}.

\textbf{Evidence anchor:} Selected full-text inspection. Corresponding 2023 arXiv manuscript: overview and architecture/skill description; bibliographic citation remains the 2024 journal record. \href{https://arxiv.org/pdf/2306.14874}{Inspected source copy}.

\paragraph{Method and contribution.} ANYmal Parkour separates perception, locomotion skills and a high-level navigation policy that selects and commands those skills. Learned scene reconstruction handles noisy or occluded observations, while a repertoire includes walking, jumping, climbing and crouching. The navigation policy reasons in relation to the capabilities available below it rather than selecting a geometric path and hoping a generic controller can realize it.

\paragraph{What the evidence establishes.} The inspected manuscript describes training modules in simulation and deploying them on ANYmal D through varied obstacle arrangements. Its platform and behavior repertoire are materially different from a G1 following planar velocity commands. The appraisal concerns the corresponding preprint architecture, not an assumption that every detail matches the final journal version.

\paragraph{Critical appraisal.} This is strong prior art against a broad claim that navigation has ignored locomotion capabilities. It does not, however, establish social acceptance or collision-free behavior among unpredictable people. Contact that is intentional and useful during parkour can be unacceptable during social passage, and high-speed terrain success is not a substitute for controlled stopping performance near a person.

\paragraph{Implication for novel work.} A response-aware G1 planner should define the executable action repertoire it assumes. Its novelty would need to explain why a calibrated low-dimensional envelope is effective or more auditable for social encounters than established capability-aware hierarchies.

\subsubsection{Mobility-aware navigation already studies delay and tracking failures}\label{paper:wheeledleg}
\textbf{Publication:} \citet{wheeledleg}. \emph{Learning robust autonomous navigation and locomotion for wheeled-legged robots}. \href{https://doi.org/10.1126/scirobotics.adi9641}{Publication record}.

\textbf{Evidence anchor:} Selected full-text inspection. Accepted manuscript: Figure 7 comparison/failure analysis, discussion and hierarchical-learning methods. \href{https://arxiv.org/pdf/2405.01792}{Inspected source copy}.

\paragraph{Method and contribution.} Lee and colleagues integrate adaptive wheeled-legged locomotion with a learned local navigation policy and a larger-scale route interface. The hierarchy lets local navigation account for what the locomotion controller can achieve. Memory supports exploration and recovery from local minima; the methods discussion describes staged training and notes that additional alternating fine-tuning gave only marginal benefits in their experiments.

\paragraph{What the evidence establishes.} The source reports kilometer-scale missions and a controlled comparison with ten trials per method in a challenging point-goal setup. Figure 7 identifies occlusion, delayed replanning, tracking overshoot and finite-memory failures rather than presenting only successful trajectories. This is directly relevant to claims about planning-execution mismatch.

\paragraph{Critical appraisal.} The paper already recognizes delayed and imperfect tracking as navigation bottlenecks, so identifying that phenomenon alone is not novel. Its morphology, mobility modes and route assumptions differ from humanoid crowd navigation. Long missions show valuable system competence but are not equivalent to an exposure-controlled contact-risk comparison across policies.

\paragraph{Implication for novel work.} A new contribution must specify a different mechanism or provide a controlled finding about calibration, uncertainty or age, beyond restating delay as a problem. The failure-analysis structure is a good model for linking an observed collision to the actual decision and execution sequence.

\subsubsection{ExBody and the upper/lower-body objective split}\label{paper:expressive}
\textbf{Publication:} \citet{expressive}. \emph{Expressive Whole-Body Control for Humanoid Robots}. \href{https://doi.org/10.15607/rss.2024.xx.107}{Publication record}.

\textbf{Evidence anchor:} Selected full-text inspection. Author preprint: retargeting, reference-state initialization, rewards, Table IV ablations and physical examples. \href{https://arxiv.org/pdf/2402.16796}{Inspected source copy}.

\paragraph{Method and contribution.} ExBody uses human motion to guide expressive upper-body behavior while relaxing lower-body imitation in favor of robust root-motion objectives. Retargeted motion supplies a structured distribution of commands and initialization states. PPO then balances expression, root movement and regularization instead of forcing the robot to reproduce every human joint trajectory despite morphological mismatch.

\paragraph{What the evidence establishes.} The platform is Unitree H1, with a 19-DoF model in the inspected version. Ablations examine reference-state initialization, random goal sampling, full-body tracking and AMP-related rewards. The discussion identifies degenerate strategies such as kneeling while satisfying parts of the objective, showing why survival or episode length alone can be misleading.

\paragraph{Critical appraisal.} Expressive movement and social navigation are connected but distinct tasks. Demonstrations of handshakes, gestures or dancing do not quantify perceived comfort during close passage. The observation/action dimensions and gains of H1 are not transferable to a G1 actor merely by copying weights or reward terms. Objective tradeoffs can also change the base response seen by a navigation planner.

\paragraph{Implication for novel work.} The paper motivates explicit arm/base coordination and posture-conditioned response tests. It also shows why a payload-policy result needs initialization and command-distribution controls before attributing improvement to a new training objective.

\subsubsection{Imitation quality and task success can move in different directions}\label{paper:deepmimic}
\textbf{Publication:} \citet{deepmimic}. \emph{DeepMimic}. \href{https://doi.org/10.1145/3197517.3201311}{Publication record}.

\textbf{Evidence anchor:} Selected full-text inspection. Author preprint: overview, Sections 10-11, Table 4 and initialization/termination ablations. \href{https://arxiv.org/pdf/1804.02717}{Inspected source copy}.

\paragraph{Method and contribution.} DeepMimic combines motion-imitation and task objectives to learn physics-based character skills, supported by reference-state initialization and termination rules that discourage invalid strategies. Policies can adapt a reference motion to a task rather than only replaying it. The approach also explores combinations of skills and morphology changes, providing foundational machinery for later humanoid motion learning.

\paragraph{What the evidence establishes.} The inspected results are simulated characters, with task-specific comparisons such as striking and throwing. Table 4 shows that the mixture of imitation and task reward changes success differently by task; imitation alone can reproduce appearance while failing the goal. The paper discloses that much of its performance evidence uses one training run because training is expensive.

\paragraph{Critical appraisal.} A natural-looking trajectory is not automatically useful, safe or hardware-realizable. Reference-phase synchronization, servo choices, reward weights and training coverage matter. The paper explicitly leaves physical robot deployment as future work in that version, so it should not be cited as direct evidence of G1 sim-to-real transfer.

\paragraph{Implication for novel work.} Use it as the foundation for imitation and initialization controls. In social navigation, report task completion, collision/fall outcomes and motion-quality proxies separately; do not let visual resemblance or a longer episode substitute for completing the mission.

\subsubsection{Dynamics randomization distributes training assumptions}\label{paper:domain}
\textbf{Publication:} \citet{domain}. \emph{Sim-to-Real Transfer of Robotic Control with Dynamics Randomization}. \href{https://doi.org/10.1109/icra.2018.8460528}{Publication record}.

\textbf{Evidence anchor:} Selected full-text inspection. Author preprint: Section IV objective, dynamics-parameter distribution and Fetch pushing-task description. \href{https://arxiv.org/pdf/1710.06537}{Inspected source copy}.

\paragraph{Method and contribution.} Dynamics randomization optimizes expected return across a distribution of simulator parameters rather than a single nominal model. A policy can use interaction history to adapt implicitly to the latent dynamics of a deployment. The inspected formulation makes the sampling distribution part of the learning problem and motivates recurrence when physical parameters are not directly observable.

\paragraph{What the evidence establishes.} The paper demonstrates the approach on a puck-pushing task with a Fetch arm, transferring a simulation-trained policy without physical-system training. This is evidence for a transfer mechanism on that task, not for humanoid walking or human-aware navigation. The physical and observation contracts differ from both RMA and the G1 payload controller.

\paragraph{Critical appraisal.} Robustness across sampled parameters does not establish coverage of arbitrary deployment dynamics. Correlations, rare extremes, delays and contact effects outside the training distribution can still matter. Randomizing a parameter also does not supply a calibrated online bound on stopping distance or body occupancy. The distinction between adaptation and certification remains essential.

\paragraph{Implication for novel work.} Use randomization as a training factor and preserve its distribution in manifests. Compare it with measured calibration and held-out perturbations, particularly for arm posture, payload inertia and stop/reverse transients that may be poorly represented by nominal locomotion training.

\subsubsection{Residual action alignment and downstream policy transfer}\label{paper:asap}
\textbf{Publication:} \citet{asap}. \emph{ASAP: Aligning Simulation and Real-World Physics for Learning Agile Humanoid Whole-Body Skills}. \href{https://arxiv.org/abs/2502.01143v3}{Publication record}.

\textbf{Evidence anchor:} Selected full-text inspection. arXiv v3: Sections III-IV, Tables III-IV, data collection and sim-to-sim/physical setup. \href{https://arxiv.org/pdf/2502.01143v3}{Inspected source copy}.

\paragraph{Method and contribution.} ASAP first learns motion tracking, collects deployment trajectories, fits a residual action model that makes simulation reproduce observed transitions, and fine-tunes the control policy in that corrected environment. The correction acts through the action interface rather than only changing the next-state prediction. This makes dynamics mismatch experimentally actionable instead of treating the simulator as a fixed ground truth.

\paragraph{What the evidence establishes.} The inspected evaluation separates open-loop transition matching from closed-loop motion imitation, compares parameter identification and residual-state alternatives, and uses cross-simulator tests before selected physical G1 tasks. It reports a 43-motion collection with difficulty levels and defines imitation success through body-tracking error, not pedestrian avoidance.

\paragraph{Critical appraisal.} Better open-loop matching need not automatically improve closed-loop behavior, which is why the paper's two-level evaluation matters. Its hardware tasks are selected to exercise transfer gaps, not a representative crowd-navigation cohort. A dynamics correction learned from those motions is not necessarily valid for loaded docking, unseen command reversals or reactive interpersonal encounters.

\paragraph{Implication for novel work.} ASAP narrows novelty around learned dynamics correction. Geoff-Sim can instead study a calibrated response envelope for a frozen policy at the navigation interface, but must compare model fidelity and downstream mission outcomes separately and demonstrate value beyond retraining or fixed conservative margins.

\subsection{Sensor-level humanoid navigation}
\subsubsection{Safety-inspired training is not an online invariance certificate}\label{paper:safehumanoid}
\textbf{Publication:} \citet{safehumanoid}. \emph{End-to-End Humanoid Robot Safe and Comfortable Locomotion Policy}. \href{https://arxiv.org/abs/2508.07611v1}{Publication record}.

\textbf{Evidence anchor:} Selected full-text inspection. arXiv v1: Sections III-V, equations 1-9, Tables I-III and physical deployment description. \href{https://arxiv.org/pdf/2508.07611v1}{Inspected source copy}.

\paragraph{Method and contribution.} Wang and colleagues map temporal LiDAR features and proprioceptive/command history to motor actions. They formulate constrained reinforcement learning with P3O, privileged training critics and barrier-inspired costs. Comfort-oriented terms shape proximity, approach behavior and smoothness. The paper explicitly says that the barrier condition is converted into a learning cost rather than used to filter actions online.

\paragraph{What the evidence establishes.} Thirty trials per simulation scenario yield rounded rates including 86\% for dynamic agents, 60\% for a narrow passage and 83\% for a suspended obstacle. The simpler P3O baseline reaches 90\% on the suspended-obstacle case. Table II reports nonzero unsafe-region time, with a 0.6 m evaluation threshold distinct from the 0.8 m training distance. G1/Livox deployment is described with onboard computation.

\paragraph{Critical appraisal.} The mathematical definition of a safe policy is stronger than what finite-weight approximate learning and these experiments establish. The descriptions of the P3O ablation leave attribution between barrier costs and comfort rewards less clean than a fully factorial design. Some reward-sign conventions merit code-level verification before reuse. Comfort is assessed by proxies rather than a reported participant study in the inspected experiments.

\paragraph{Implication for novel work.} This is direct prior art for sensor-driven G1 avoidance. A modular method must compete on matched sensing, timing, contacts and liveness, while stating whether its contribution is interpretability, calibrated uncertainty, compute efficiency or a stronger justified guarantee.

\subsubsection{Omni-Perception and the metric ambiguity of success plus collision}\label{paper:omni}
\textbf{Publication:} \citet{omni}. \emph{Omni-Perception: Omnidirectional Collision Avoidance for Legged Locomotion in Dynamic Environments}. \href{https://arxiv.org/abs/2505.19214}{Publication record}.

\textbf{Evidence anchor:} Selected full-text inspection. Retrieved arXiv v2: Section 3 architecture; Sections 4-6; Tables 1, 3 and 4; failure cases. \href{https://arxiv.org/pdf/2505.19214}{Inspected source copy}.

\paragraph{Method and contribution.} Omni-Perception processes LiDAR history through proximal and distal pathways: near geometry receives more detailed sampling, while farther geometry is compressed. Recurrent features feed a locomotion policy, and a cross-platform raycasting toolkit supports training with sensor-pattern and noise variation. Avoidance rewards include sector-based directional behavior and clearance-related terms alongside commanded motion.

\paragraph{What the evidence establishes.} The paper separates LiDAR rendering throughput from policy ablations and physical scenarios. Table 3 uses six dynamic obstacles and thirty trials, reporting 76.7\% success and 56.7\% collision for the proposed method. Table 4 includes 27/30 successful moving-human trials, but a small-obstacle case favors the native Unitree comparator. The hardware evidence spans legged contexts and must not be assumed to be a matched G1-only benchmark.

\paragraph{Critical appraisal.} Success and collision entries in Table 3 are not complementary, and the inspected reporting does not justify converting them into a collision-free completion rate. Similar-looking point clouds are useful qualitative sensor evidence but not a calibrated sensor-error distribution. The limitations identify lost thin-object detail from downsampling and failures in dense vegetation, exposing a perception-compression tradeoff.

\paragraph{Implication for novel work.} This paper rules out broad novelty claims about raw-LiDAR end-to-end legged avoidance or multi-simulator LiDAR rendering. A useful project comparison should separate sensor fidelity, runtime age, obstacle recall, goal completion and actual contact, with transparent event definitions.

\subsubsection{Generative navigation, dynamic adaptation and cross-embodiment evidence}\label{paper:flux}
\textbf{Publication:} \citet{flux}. \emph{FLUX: Accelerating Cross-Embodiment Generative Navigation Policies via Rectified Flow and Static-to-Dynamic Learning}. \href{https://arxiv.org/abs/2603.12806v1}{Publication record}.

\textbf{Evidence anchor:} Selected full-text inspection. arXiv v1: dynamic-scene protocol, metric definitions, Sections V-VI and Tables II-IV. \href{https://arxiv.org/pdf/2603.12806v1}{Inspected source copy}.

\paragraph{Method and contribution.} FLUX uses rectified-flow trajectory generation and a static-to-dynamic learning curriculum. Visual history and depth condition a transformer-based policy, with dynamic-task refinement after acquiring static geometric priors. DynBench introduces continuing pedestrian activity and structured routes, so the challenge is more than avoiding agents that stop permanently at their destinations.

\paragraph{What the evidence establishes.} Table II reports SocialNav success of 64.0 versus 59.3 for NavDP in that protocol. Physical demonstrations include wheeled, quadrupedal and G1 platforms, with RGB-D sensing and remote computation described in the paper. The reported 216.5 ms selected inference configuration uses a different timing boundary and hardware from the project's local planner. These values should not be placed in a common leaderboard.

\paragraph{Critical appraisal.} The metric named Collision Count is defined through distance below 0.45 m, not direct contact. The social-distance measure also needs its direction-of-improvement read carefully: time inside a proximity region is not a positive comfort score merely because the metric is named compliance. Short inference improvements do not by themselves quantify acquisition-to-execution freshness or human comfort.

\paragraph{Implication for novel work.} FLUX establishes close prior art for cross-embodiment dynamic navigation and G1 transfer. The remaining opening is a precise response/observation contract and auditable downstream evidence, not the first combination of a humanoid and a generative or sensor-driven navigation pipeline.

\subsection{Enabling technology and systems}
\subsubsection{Physics simulation enables experiments but does not validate their model}\label{paper:mujoco}
\textbf{Publication:} \citet{mujoco}. \emph{MuJoCo: A physics engine for model-based control}. \href{https://doi.org/10.1109/iros.2012.6386109}{Publication record}.

\textbf{Evidence anchor:} Record-limited inspection. Verified 2012 DOI identity; accessible current official overview for software context. Attempts to obtain the original author PDF were unsuccessful.

\paragraph{Method and contribution.} MuJoCo is the physics foundation for the project, supporting articulated bodies, contacts, actuation and sensor-related calculations. The current official overview explains its model/data separation and articulated simulation approach. This appraisal distinguishes that accessible software context from the original 2012 paper, whose full experimental text was not recovered during this review.

\paragraph{What the evidence establishes.} The primary publication exists and is bibliographically verified. The project can run native simulator tests and compare model behavior, but those checks are evidence about its implementation, not independent validation of every contact or sensor assumption. No numerical claim from the inaccessible original paper is newly attributed here.

\paragraph{Critical appraisal.} A physically simulated result remains conditional on geometry, inertial parameters, actuator gains, contact settings, integration and sensor timing. Visual realism, a named physics engine and a high simulation rate answer different questions. Likewise, a magnetic attachment is an explicit development mechanism, not evidence of finger-force closure simply because the rest of the scene obeys physics.

\paragraph{Implication for novel work.} Record the actual simulator version and model/patch identity, calibrate critical behavior and separate contact semantics from visual overlays. Simulator integration is enabling engineering; a scientific contribution requires a new question or validated method tested with it.

\subsubsection{Middleware timing must be measured at the right boundary}\label{paper:ros2}
\textbf{Publication:} \citet{ros2}. \emph{Robot Operating System 2: Design, architecture, and uses in the wild}. \href{https://doi.org/10.1126/scirobotics.abm6074}{Publication record}.

\textbf{Evidence anchor:} Selected full-text inspection. Author preprint: architecture overview, networking experiments and deployment-case discussion. \href{https://arxiv.org/pdf/2211.07752}{Inspected source copy}.

\paragraph{Method and contribution.} ROS 2 supplies modular communication, DDS-based middleware abstraction and lifecycle concepts for composing robotics systems. Its design supports different transport and execution choices rather than prescribing a single timing behavior. The paper connects those architectural decisions to practical systems and contrasts intra-process, composed and multiprocess communication.

\paragraph{What the evidence establishes.} The inspected networking experiments specify hardware, message sizes, publication rates, middleware and packet-loss conditions. Sub-millisecond p95 results in some intra-process conditions are meaningful at that measured boundary. They do not include an arbitrary robot's detector, planner, policy inference, bridge and physical response.

\paragraph{Critical appraisal.} Middleware reliability and physical safety are different claims. A successfully delivered message may contain an old observation; a fast callback may wait behind other work; a published target is not a motor acknowledgement. Quality-of-service choices can trade loss, buffering and freshness, so aggregate node timing cannot reconstruct the full causal pipeline without timestamps and identity.

\paragraph{Implication for novel work.} Instrument acquisition, decision, command consumption and response separately. A source-audited timing study can be scientifically useful if it tests a new causal hypothesis, while simply using ROS 2 or reporting average callback time remains infrastructure work.

\subsubsection{Behavior-tree orchestration and long-run field evidence}\label{paper:nav2}
\textbf{Publication:} \citet{nav2}. \emph{The Marathon 2: A Navigation System}. \href{https://doi.org/10.1109/iros45743.2020.9341207}{Publication record}.

\textbf{Evidence anchor:} Selected full-text inspection. Author preprint: architecture overview, experimental setup and Tables II-III. \href{https://arxiv.org/pdf/2003.00368}{Inspected source copy}.

\paragraph{Method and contribution.} The Marathon 2 presents Navigation2's modular navigation architecture and behavior-tree task orchestration. Planning, control and recovery are composed as interacting services/actions rather than a single monolithic navigation function. That structure makes it possible to adapt the task logic and plugins, but it also creates interfaces whose readiness and timing matter.

\paragraph{What the evidence establishes.} The campus experiment uses RB1 and TIAGo platforms, reporting 37.4 combined miles, 22.8 hours and 168 recoveries, with zero recorded collisions and emergency stops in that study. Both robots are limited to 0.45 m/s; one listed configuration uses external navigation computation and another onboard computation. These are important conditions of the field evidence.

\paragraph{Critical appraisal.} Long-duration operation is stronger evidence of system integration than a short demo, but it does not establish a humanoid control model or a universal collision probability. Recovery count also matters: successful operation can involve repeated interventions by the autonomy stack without human intervention. Task success should not erase the cost or cause of those recoveries.

\paragraph{Implication for novel work.} Treat Nav2 as established infrastructure and a useful source of lifecycle/recovery patterns. Research novelty should concern how measured G1 response, observation freshness or planner-supervisor consistency changes outcomes, not the presence of a behavior tree or repeated waypoint goals.

\subsubsection{Mapping modes trade completeness against freshness}\label{paper:slam}
\textbf{Publication:} \citet{slam}. \emph{SLAM Toolbox: SLAM for the dynamic world}. \href{https://doi.org/10.21105/joss.02783}{Publication record}.

\textbf{Evidence anchor:} Selected full-text inspection. JOSS paper: summary, feature description, graph serialization and synchronous/asynchronous mode discussion. \href{https://www.theoj.org/joss-papers/joss.02783/10.21105.joss.02783.pdf}{Inspected source copy}.

\paragraph{Method and contribution.} SLAM Toolbox provides graph-based mapping, localization and persistent map/session operations. The paper highlights serialization of raw data and the pose graph, allowing a session to be resumed or refined. Synchronous operation can queue observations for mapping quality, whereas asynchronous processing favors current localization by admitting new work when the previous update completes.

\paragraph{What the evidence establishes.} The inspected feature discussion explains practical mapping and map-maintenance capabilities. It is a software paper, not a randomized controlled comparison of social navigation policies. Its treatment of large and changing environments should therefore be understood as mapping functionality and application experience rather than direct evidence of dynamic-person avoidance.

\paragraph{Critical appraisal.} The distinction between map completeness and freshness is central to downstream planning. A good eventual map can coexist with decisions made from old state, while dropping scans can improve freshness but lose geometry. Localization corrections can also move goals and obstacle representations unless the whole pipeline uses consistent frame/time semantics.

\paragraph{Implication for novel work.} Record which mapping mode and update policy ran. Test whether map quality, pose corrections or delivered-route age actually explain mission failures instead of assuming that successful SLAM implies reliable navigation through the crowd.
